% sec-10: how the document was built, and what survives it.  Figure 20.
\section{Build Method and Reproducibility}

\subsection{The Eight-Stage Build}

\draftinstr{One paragraph plus a compact eight-row table: the five diagram
stages and the three paper stages, each with its output directory, its figure
numbers and its commit count. Source:
\appfile{funding/capitalization-plan/sub-prompts/README.md}. State the two rules
that shaped the build: one commit per file the moment it is written, so the
branch can be monitored without intervention, and no `.py` file anywhere, so the
three \texttt{lint-and-format} jobs stay green.}

\subsection{Why the Diagram Split Is Uneven}

\draftinstr{One paragraph. Five mermaid, three plantuml, five d2, three
diagrams, four graphviz. The split follows purpose and not quota: the two
five-counts fall where the paper argues in sequence and where it argues in
tables, which is what a capitalization plan does most. Source:
\appfile{funding/capitalization-plan/sub-prompts/README.md}. Do not restate the
per-figure reasoning, which is in the five directory READMEs.}

\subsection{What Survives a Stop}

\figslot{diagrams-python-type / custody clusters}{9.0}{Four custodians and what
each still holds if the programme stops.\\Everything a third party needs to
reproduce the computation sits\\with the first custodian, which needs no one to
stay in business.}

\draftinstr{After Figure 20, one paragraph on the reproducibility count: five of
seven work products are fully reproducible from the public custodian alone, and
the two that are not require source documents that never leave the site. Then
the argument this figure exists to make, stated once and not repeated: if the
company dissolves at any milestone, the computational work does not become
unavailable, because it was never in the company's custody in the first place.
Source: \appfile{funding/capitalization-plan/diagrams-python/fig-20-artifact-custody.md}.
Cite \texttt{ecfr2026part312records} and \texttt{repo2026}.}

\subsection{Verification of This Document}

\draftinstr{A short compile record in prose, not a table: pdfLaTeX twice plus
BibTeX, zero errors, zero overfull boxes, zero undefined citations, zero
undefined references, no raster image anywhere, and the figure-to-caption
distance identical at 6.06 pt for every one of the twenty figures. Name the
audit that proves the last claim, which is that the count of
\texttt{\textbackslash vspace\{-0.65cm\}} equals the count of
\texttt{\textbackslash figcaption} plus the count of \texttt{\textbackslash
tabcap}. Then name the three defects found and fixed during the build, drawn
from the stage READMEs, so the record is a record and not an assertion.}
